\documentclass[12pt,a4paper]{article}
\usepackage[a4paper,left=30mm,right=15mm,top=20mm,bottom=20mm,footskip=10mm]{geometry}
\usepackage{fontspec}
\usepackage{polyglossia}
\setmainlanguage{russian}
\setotherlanguage{english}
\IfFontExistsTF{Times New Roman}{
  \setmainfont{Times New Roman}
}{
  \IfFontExistsTF{Liberation Serif}{
    \setmainfont{Liberation Serif}
  }{
    \setmainfont{TeX Gyre Termes}
  }
}
\IfFontExistsTF{Arial}{
  \setsansfont{Arial}
}{
  \IfFontExistsTF{DejaVu Sans}{
    \setsansfont{DejaVu Sans}
  }{
    \setsansfont{TeX Gyre Heros}
  }
}
\IfFontExistsTF{DejaVu Sans Mono}{
  \setmonofont{DejaVu Sans Mono}
}{
  \IfFontExistsTF{Courier New}{
    \setmonofont{Courier New}
  }{}
}
\usepackage{amsmath,amssymb,booktabs,tabularx,array,longtable}
\usepackage{graphicx,caption,float,xcolor,microtype}
\usepackage[hidelinks]{hyperref}
\usepackage{enumitem}
\usepackage{setspace}

\graphicspath{{report/figures/}}
\input{report/generated_metrics.tex}
\hypersetup{colorlinks=false}
\pagestyle{plain}
\onehalfspacing
\setlength{\parindent}{1.25cm}
\setlength{\parskip}{0pt}
\setlength{\emergencystretch}{3em}
\raggedbottom
\setlist{leftmargin=1.25cm,itemsep=0pt,topsep=3pt,parsep=0pt}
\DeclareCaptionLabelSeparator{gosthyphen}{\space-\space}
\captionsetup[figure]{name=Рисунок,labelsep=gosthyphen,font=normalsize,
  labelfont=bf,justification=centering,singlelinecheck=false}
\captionsetup[table]{name=Таблица,labelsep=gosthyphen,font=normalsize,
  labelfont=bf,justification=raggedright,singlelinecheck=false,position=top}
\renewcommand{\contentsname}{\centering СОДЕРЖАНИЕ}
\renewcommand{\refname}{\centering СПИСОК ИСПОЛЬЗОВАННЫХ ИСТОЧНИКОВ}
\makeatletter
\renewcommand\section{\@startsection{section}{1}{0pt}
  {-2.8ex plus -0.8ex minus -0.2ex}{1.2ex plus 0.2ex}
  {\centering\normalfont\normalsize\bfseries}}
\renewcommand\subsection{\@startsection{subsection}{2}{0pt}
  {-2.2ex plus -0.6ex minus -0.2ex}{0.8ex plus 0.2ex}
  {\normalfont\normalsize\bfseries}}
\renewcommand{\@seccntformat}[1]{\csname the#1\endcsname\hspace{0.5em}}
\makeatother
\AtBeginDocument{%
  \renewcommand{\contentsname}{\centering СОДЕРЖАНИЕ}%
  \renewcommand{\refname}{\centering СПИСОК ИСПОЛЬЗОВАННЫХ ИСТОЧНИКОВ}%
}

\newcommand{\code}[1]{\texttt{#1}}
\newcommand{\term}[1]{\textit{#1}}

\begin{document}
\begin{titlepage}
\thispagestyle{empty}
\begin{center}
\normalsize
Министерство науки и высшего образования Российской Федерации\\
ФЕДЕРАЛЬНОЕ ГОСУДАРСТВЕННОЕ АВТОНОМНОЕ\\
ОБРАЗОВАТЕЛЬНОЕ УЧРЕЖДЕНИЕ ВЫСШЕГО ОБРАЗОВАНИЯ\\[2mm]
\textbf{НАЦИОНАЛЬНЫЙ ИССЛЕДОВАТЕЛЬСКИЙ УНИВЕРСИТЕТ ИТМО}\\
\textbf{ИТМО} ITMO University\\[7mm]
\textbf{Мегафакультет}\\
Мегафакультет компьютерных технологий и управления\\
\textbf{Факультет}\\
Институт математики\\[7mm]
Образовательная программа: «Математические технологии системного анализа» 2025\\[3mm]
Направление подготовки: 27.04.03 - Системный анализ и управление
\end{center}

\vspace{13mm}
\begin{center}
{\fontsize{16}{19}\selectfont О Т Ч Е Т}\\[3mm]
о \textit{научно-исследовательской} практике\\[8mm]
\textbf{Тема задания:} Восстановление локальной стохастической динамики\\
для раннего обнаружения критических переходов
\end{center}

\vspace{9mm}
\begin{flushleft}
\textbf{Обучающийся:} Д.~А.~Ластовецкий, № группы R4160\\[3mm]
\textbf{Уровень образования:} магистратура\\[3mm]
\textbf{Направление подготовки:} 27.04.03 Системный анализ и управление\\[3mm]
\textbf{Образовательная программа:} «Математические технологии системного анализа» 2025\\[3mm]
\textbf{Руководитель практики от университета ИТМО:}\\
Леоненко Василий Николаевич, кандидат физико-математических наук, доцент
\end{flushleft}

\vspace{6mm}
\begin{flushright}
\textbf{Дата:} 16.09.2026
\end{flushright}

\vfill
\begin{center}
Санкт-Петербург\\
2026 г.
\end{center}
\end{titlepage}
\setcounter{page}{2}

\setcounter{tocdepth}{2}
\begingroup
\singlespacing
\tableofcontents
\endgroup
\clearpage

\section*{Введение}
\addcontentsline{toc}{section}{Введение}

\textbf{Актуальность исследования.} Сложные системы способны долго оставаться в относительно устойчивом режиме,
а затем быстро переходить в качественно иное состояние. Для эпидемиологии это
может быть начало волны заболеваемости; для экологии --- смена устойчивого
состояния экосистемы; для финансов --- переход к режиму повышенной
волатильности. Наблюдателю доступен, как правило, не полный набор переменных
системы, а один или несколько зашумленных временных рядов. Кроме того,
наблюдения могут поступать редко, содержать пропуски и отражать одновременно
динамику объекта и ошибки измерения.

Под \term{критическим переходом} далее понимается резкое качественное
изменение режима системы при приближении управляющего параметра к пороговому
значению. Вблизи такого порога малое изменение условий может привести к
переходу в другое устойчивое состояние~\cite{scheffer2009}.

Если переходу предшествует потеря устойчивости текущего состояния, система
после малого возмущения возвращается к равновесию все медленнее. Этот эффект
называется \term{критическим замедлением}. Его операциональный признак ---
увеличение времени восстановления после малого возмущения; в данных он часто
проявляется ростом дисперсии и лаговой автокорреляции~\cite{scheffer2009}.
Статистики, которые меняются до перехода и могут быть вычислены по наблюдениям,
называются \term{сигналами раннего предупреждения} (early warning signals,
EWS). Такие статистики удобны, поскольку почти не
требуют знания механизма. Однако один и тот же эмпирический паттерн может быть
вызван разными причинами. Например, дисперсия увеличивается как при ослаблении
восстанавливающей силы, так и при усилении внешнего шума. Для практического
решения эти случаи принципиально различны: первый поддерживает гипотезу о
приближении потери устойчивости, второй --- гипотезу об усилении возмущений
при сохранении локальной динамики.

Отсюда возникает задача перейти от описания наблюдаемого ряда к
идентификации механизма. Нужен метод, который в каждом локальном окне оценивает
скорость возврата к равновесию, интенсивность процессного шума и, по
возможности, ошибку наблюдения. Такой переход делает сигнал более
интерпретируемым, но одновременно усложняет оценивание: короткое окно содержит
мало информации, а гибкая модель легко принимает шум измерения за нелинейную
динамику.

\textbf{Научная проблема.} Рассматривается стохастическая динамическая система
\begin{equation}
\label{eq:sde}
 dX_t=f(X_t,\theta_t)\,dt+G(X_t,\theta_t)\,dW_t,
\end{equation}

где $X_t$ --- состояние системы в момент $t$; $f$ --- детерминированный снос,
задающий среднюю локальную тенденцию; $G$ --- коэффициент диффузии,
определяющий интенсивность случайных возмущений; $W_t$ --- винеровский
процесс; $\theta_t$ --- медленно меняющийся параметр. Стохастическим
дифференциальным уравнением (СДУ) называется уравнение вида~\eqref{eq:sde}, в
котором эволюция состояния включает детерминированную и случайную компоненты;
раздельное восстановление $f$ и $G$ по дискретным наблюдениям составляет
задачу идентификации СДУ~\cite{bonalli2026}.

\textbf{Исследовательский вопрос.} Исследование отвечает на вопрос:

\begin{quote}
Можно ли по коротким, шумным и неполным временным рядам достаточно надежно
восстанавливать изменение локальной устойчивости и использовать его для
раннего обнаружения критического перехода, одновременно отличая потерю
устойчивости от простого роста случайных возмущений?
\end{quote}

Вопрос затрагивает три взаимосвязанных аспекта. Под \term{идентифицируемостью}
понимается возможность различить значения параметров по распределению
наблюдаемых данных; здесь это означает раздельное оценивание сноса и диффузии.
\term{Робастность} характеризует сохранение основного вывода при изменении
сетки наблюдений, уровня шума измерения и доли пропусков. Под
\term{калибровкой неопределенности} понимается соответствие заявленной
вероятности интервала фактической частоте включения истинного значения.

\textbf{Объект исследования.} Временные ряды стохастических динамических
систем, локальная устойчивость которых может меняться перед переходом между
режимами.

\textbf{Предмет исследования.} Методы локального восстановления сноса,
диффузии и скорости возврата к равновесию, а также правила формирования на их
основе ранней тревоги.

\textbf{Цель работы.} Определить условия, при которых восстановление
локальной стохастической динамики дает более специфичный, устойчивый и
интерпретируемый сигнал критического перехода, чем классические статистические
индикаторы и непосредственно обучаемые модели машинного обучения.

\textbf{Задачи работы.} Для достижения цели решаются следующие задачи:

\begin{enumerate}
\item формализовать связь между локальной устойчивостью, диффузией,
наблюдаемой дисперсией и критическим замедлением;
\item построить контролируемые синтетические сценарии потери устойчивости,
роста шума и стационарного поведения с известными истинными параметрами;
\item включить как линейный OU-контроль, так и нелинейные нормальные формы
седло-узел, транскритическую и бистабильную с диффузией, зависящей от состояния;
\item реализовать единый набор классических, параметрических, разреженных,
байесовских и нейросетевых моделей и сравнить их на одинаковых окнах;
\item оценить не только качество классификации, но и восстановление параметров,
долю успешных подгонок, покрытие интервалов, ложные тревоги и время упреждения;
\item проверить устойчивость выводов к разрежению наблюдений, ошибке измерения
и случайным пропускам;
\item воспроизвести критерий начала эпидемической волны и сравнить
классификаторы на публичном недельном ряду ОРВИ;
\item сопоставить локальную SDE-модель с Persistence, LSTM и GRU в прогнозе
COVID-индикаторов со скользящим началом (rolling origin) на одинаковых
горизонтах прогнозирования.
\end{enumerate}

\textbf{Проверяемые гипотезы. H1. Специфичность механизма.} Совместная оценка $\kappa$ и $\sigma$
должна реже смешивать потерю устойчивости с ростом шума, чем одна дисперсия.

\textbf{H2. Цена гибкости.} SINDy и Neural SDE способны лучше описывать
нелинейный снос, но на коротких окнах их преимущество может исчезать из-за
высокой дисперсии оценки и технических отказов.

\textbf{H3. Роль модели наблюдений.} Явное разделение процессного шума и
шума измерения должно повышать долю успешных оценок и устойчивость на неполных
данных.

\textbf{H4. Ограниченная переносимость.} Модель, обученная или настроенная на
синтетических переходах, не обязана сохранять калибровку на реальном
эпидемиологическом ряду без адаптации домена.

\textbf{Научная новизна.} Модели разных классов сопоставляются не только по
ROC-AUC, но и по восстановлению физически интерпретируемых параметров и
последовательным тревогам. Полное апостериорное распределение
spike-and-slab усредняет восемь моделей кубической библиотеки и переносит
неопределенность выбора структуры на интервалы $\kappa$. Neural SDE со скрытым
состоянием отделяет латентную траекторию от наблюдений и тем самым проверяет,
помогает ли явная модель наблюдений в условиях пропусков. Все отрицательные
результаты и неуспешные попытки оценивания сохраняются, а не исключаются из
итоговой таблицы.

\textbf{Практическая значимость.} Результатом работы стала воспроизводимая схема
эпидемиологического анализа. Она требует одинаковых горизонтов, хронологического
разбиения, калибровки порога только на обучающих траекториях и совместного
отчета о точности, доле успешных оценок, покрытии интервалов и ложных тревогах.
\textbf{Границы исследования.} Работа не претендует на готовую систему
медицинского оповещения: рассмотрены одномерные модели и два конкретных
набора данных, а эпидемиологические выводы являются ретроспективными.

\textbf{Структура отчета.} Раздел~1 объясняет, почему классические EWS неоднозначны и как локальная SDE
делает механизм явным. Раздел~2 фиксирует данные, сценарии и правила сравнения.
Раздел~3 описывает модели от простых базовых методов до Bayesian SINDy и
Neural SDE со скрытым состоянием. В разделах~4 и~5 приведены синтетические и
эпидемиологические результаты.
Раздел~6 переводит метрики в содержательные выводы и формулирует следующий
этап. Полные таблицы вынесены в приложение, чтобы основной текст оставался
читаемым.

\textbf{Терминологические соглашения.} В основном тексте используются русские термины: \term{снос},
\term{диффузия}, \term{процессный шум}, \term{шум измерения},
\term{модель наблюдений}, \term{доля успешных оценок} и
\term{покрытие интервала}. Английские названия сохранены только для
общепринятых метрик и собственных названий алгоритмов, например ROC-AUC,
SINDy, Neural SDE, Easy Ensemble и rolling origin. Такое соглашение устраняет
смешение языков, не затрудняя сопоставление с литературой и кодом.

\section{Теоретическая рамка}

\subsection{Основные понятия локальной динамики}

Точка $x^*$ называется \term{равновесием}, если $f(x^*)=0$. В одномерной
системе равновесие локально устойчиво, когда $f'(x^*)<0$: после малого
отклонения снос направляет траекторию обратно к $x^*$. В работе скорость
локального восстановления обозначается
$\kappa=-f'(x^*)>0$. Чем меньше $\kappa$, тем слабее восстанавливающее
воздействие и тем дольше система возвращается к равновесию. Такое определение
соответствует стандартной локальной классификации одномерных
равновесий~\cite{strogatz2015}.

Под \term{процессным шумом} понимаются случайные возмущения самой динамики,
представленные диффузионным членом $G\,dW_t$ в уравнении~\eqref{eq:sde}.
\term{Шум измерения} возникает при регистрации состояния и задается отдельно,
например уравнением $Y_t=X_t+\varepsilon_t$. Совместное задание уравнения
состояния и условного распределения наблюдений образует \term{модель в
пространстве состояний}. Разделение предотвращает отнесение ошибки измерения к
процессной диффузии и соответствует постановке последовательного
оценивания скрытого состояния по шумным наблюдениям~\cite{kalman1960}.

\subsection{Критическое замедление и неоднозначность дисперсии}

В малой окрестности устойчивого равновесия $x_t^*$ одномерную систему можно
локально представить процессом Орнштейна--Уленбека:
\begin{equation}
\label{eq:ou}
 dX_t=-\kappa_t(X_t-x_t^*)\,dt+\sigma_t\,dW_t,
 \qquad
 \operatorname{Var}(X_t)\approx\frac{\sigma_t^2}{2\kappa_t}.
\end{equation}
Процесс Орнштейна--Уленбека --- линейный гауссовский марковский процесс с
возвратом к среднему; здесь он используется как локальная аппроксимация
динамики около устойчивого равновесия~\cite{carpenter2011}.
Параметр $\kappa_t=-\partial f/\partial x|_{x_t^*}>0$ измеряет скорость
восстановления: чем он меньше, тем медленнее затухает отклонение. При потере
устойчивости $\kappa_t\to0$, поэтому растут дисперсия, AR(1) и характерное время
возврата. Однако та же формула показывает ключевую проблему: дисперсия растет и
при $\sigma_t\uparrow$ с неизменным $\kappa_t$. Наблюдаемая дисперсия является
отношением двух механизмов и сама по себе их не идентифицирует.

Классические EWS остаются важными базовыми методами и полезны при слабом знании
модели. Непараметрическое восстановление функции сноса может извлекать
дополнительную структуру из неизвестной системы~\cite{carpenter2011}, а
обобщенное моделирование занимает промежуточное положение между
безмодельной статистикой и полностью заданной механистической
моделью~\cite{lade2012}. Глубокое обучение расширяет этот подход, обучаясь на
множестве нормальных форм переходов~\cite{bury2021}. В то же время современные
результаты по непараметрическому обучению СДУ подчеркивают, что скорость
восстановления сноса и диффузии зависит от регулярности и объема
наблюдений~\cite{bonalli2026}.

В работе разделяются три разных утверждения о природе ранних сигналов:

\begin{enumerate}
\item \textbf{Ретроспективное различение:} позднее окно сценария потери
устойчивости получает более высокий сигнал, чем позднее окно сценария роста
шума. Качество различения измеряется ROC-AUC.
\item \textbf{Правильная идентификация:} оценки $\widehat\kappa$,
$\widehat\sigma$ и $\widehat\tau$ близки к известным истинным значениям. Ее
измеряют RMSE и эмпирическое покрытие интервалов.
\item \textbf{Причинная ранняя тревога:} сигнал пересекает заранее
откалиброванный порог до момента перехода. Здесь важны доля обнаруженных
переходов, частота ложных тревог и время упреждения.
\end{enumerate}

Высокий AUC не гарантирует раннее обнаружение. Модель может хорошо ранжировать
две заранее выбранные точки, но не пересечь строгий последовательный порог.
Точно так же хороший AUC при низкой доле успешных оценок описывает только
удобное подмножество окон. Поэтому все метрики рассматриваются совместно.

\section{Данные и схема эксперимента}

\subsection{Синтетические сценарии}

Для каждой системы создаются три сценария. В сценарии \term{потери устойчивости}
$\kappa(t)$ плавно уменьшается от 0.8 до 0.1 к $t_c=80$, а процессная диффузия
остается постоянной. В сценарии \term{роста шума} локальная устойчивость
фиксирована, но $\sigma(t)$ растет так, чтобы локальная OU-дисперсия совпадала
со сценарием потери устойчивости. В \term{стационарном сценарии} оба параметра
постоянны; он служит для калибровки последовательного порога.

Помимо линейной OU-модели использованы три нелинейные нормальные формы.
\term{Нормальная форма} --- минимальная локальная модель, сохраняющая тип
бифуркации после допустимых замен переменных и параметров; седло-узел и
транскритическая бифуркация относятся к стандартным одномерным формам
потери или обмена устойчивостью~\cite{strogatz2015}:
\begin{equation}
\label{eq:normal-forms}
\begin{array}{lll}
\text{седло-узел:} & f(x)=r-x^2, & x^*=\sqrt r,\quad\kappa=2\sqrt r,\\
\text{транскритическая:} & f(x)=rx-x^2, & x^*=r,\quad\kappa=r,\\
\text{бистабильная (сборка):} & f(x)=r+x-x^3, & \kappa=3(x^*)^2-1.
\end{array}
\end{equation}
Управляющий параметр выбирается так, чтобы истинная локальная $\kappa(t)$
имела один и тот же график, несмотря на разную глобальную геометрию. Диффузия
зависит от состояния:
\begin{equation}
\label{eq:state-dependent-diffusion}
 g(x,t)=\sigma(t)\{1+0.3\tanh[x-x^*(t)]\}.
\end{equation}
Такая конструкция проверяет не запоминание OU-геометрии, а перенос метода на
разные формы сноса.

\subsection{Режимы наблюдений}

Основные параметры базового и стрессового режимов приведены в
таблице~\ref{tab:observation-regimes}.

\begin{table}[htbp]
\caption{Параметры режимов наблюдений}
\label{tab:observation-regimes}
\centering\small
\begin{tabularx}{\textwidth}{>{\bfseries}p{39mm}XX}
\toprule
Компонент & Базовый режим & Стрессовый режим \\
\midrule
Сетка наблюдений & частая, шаг 0.2 & редкая, шаг 0.8 \\
Шум измерения & малый & повышенный \\
Пропуски & отсутствуют & 30\% MCAR \\
Окно & ширина 16 & ширина 16 \\
Контрольные точки & $t=25$ и $t=78$ & $t=25$ и $t=78$ \\
\bottomrule
\end{tabularx}
\end{table}

Сокращение MCAR означает полностью случайный механизм пропусков: вероятность
отсутствия значения не зависит ни от наблюдаемых, ни от пропущенных
величин~\cite{rubin1976}.

В OU-эксперименте использованы 36 реплик на сценарий: реплики 0--23 образуют
обучающую выборку, 24--35 --- отложенную тестовую выборку. В нелинейном
эксперименте использованы 30 реплик, разделенные как 0--19 и 20--29. Порог
тревоги равен 95-му перцентилю максимального сигнала вдоль каждой стационарной
обучающей траектории. Такая калибровка учитывает многократную проверку порога и
не использует будущую тестовую информацию.

\subsection{Эпидемиологические данные}

Первый прикладной набор --- 1543 недельных наблюдения заболеваемости ОРВИ в
Санкт-Петербурге с официальным бинарным маркером эпидемического периода.
Публичный ряд охватывает 1990-01-01--2019-07-22. Началом волны считается
переход от неэпидемического режима к ближайшему размеченному эпидемическому
периоду в пределах горизонта $h=4$ или $5$ недель. Признаки строятся только из
предшествующих значений: асимметрия на окне 20 недель и дисперсия на окнах 15,
20 и 30 недель. Такое определение буквально воспроизводит критерий статьи ~\cite{lazutov2026}.

Второй набор --- ежедневные COVID-19 показатели Санкт-Петербурга из открытого
репозитория мониторинга. Целями служат 14-дневные скользящие среднее и
дисперсия доли положительных тестов; горизонты равны 14 и 28 дням.

\subsection{Как читать метрики}

Определения и направление улучшения метрик сведены в
таблицу~\ref{tab:metrics}.

\begin{table}[htbp]
\caption{Метрики качества и их содержательная интерпретация}
\label{tab:metrics}
\centering\small
\begin{tabularx}{\textwidth}{>{\bfseries}p{31mm}p{32mm}X}
\toprule
Метрика & Направление оптимального значения & Содержательный смысл \\
\midrule
ROC-AUC / PR-AUC & выше & качество ранжирования; PR-AUC предпочтительнее при редком начале волны~\cite{saito2015} \\
RMSE параметра & ниже & близость оценки к истинному параметру синтетической системы \\
Доля успешных оценок & ближе к 100\% & доля окон, где метод вообще вернул допустимую оценку \\
Покрытие интервала & около 95\% & частота включения истины в заявленный 95\% интервал \\
Доля обнаружений & выше & доля переходов, обнаруженных до контрольного момента \\
Частота ложных тревог & ниже & доля тревог при росте шума или стационарности \\
Время упреждения & выше при контроле ложных тревог & фактическое упреждение только среди обнаруженных путей \\
NRMSE & ниже & ошибка прогноза, нормированная на вариативность цели \\
\bottomrule
\end{tabularx}
\end{table}

\section{Сравниваемые модели}

\subsection{Классические и параметрические базовые модели}

Классическая группа включает дисперсию, AR(1), скорость восстановления,
коэффициент вариации, асимметрию и их ансамбль. Moment OU получает
$\widehat\kappa$ из локальной AR(1)-аппроксимации и оценивает диффузию по
стационарной дисперсии.
Kalman OU рассматривает истинное состояние как скрытый процесс, использует
точную переходную плотность OU-процесса для неравномерных интервалов и
совместно максимизирует правдоподобие по $\kappa$, $\sigma$ и стандартному
отклонению шума измерения $\tau$. Эта модель важна как простой, но структурно
корректный базовый метод. Фильтр Калмана --- рекурсивная процедура оценивания
скрытого состояния линейной гауссовской системы по последовательности шумных
наблюдений; она обновляет оценку и ее ковариацию при поступлении каждого нового
измерения~\cite{kalman1960}.

\subsection{SINDy и Bayesian SINDy}

Метод разреженной идентификации нелинейной динамики (Sparse Identification of
Nonlinear Dynamics, SINDy) представляет неизвестную правую часть уравнения как
разреженную линейную комбинацию заранее заданных функций-кандидатов и оценивает
небольшое число ненулевых коэффициентов по данным~\cite{brunton2016}.
В данной работе SINDy-STLSQ выбирает члены кубической библиотеки
$\Theta(x)=(1,x,x^2,x^3)$ последовательной пороговой гребневой регрессией.
Для нелинейной
модели сначала находится устойчивый корень $\widehat x^*$, а затем
\begin{equation}
\label{eq:kappa-estimate}
 \widehat\kappa=-\widehat f'(\widehat x^*).
\end{equation}

\term{Bayesian SINDy} заменяет единственную разреженную оценку апостериорным
распределением по коэффициентам и структурам модели. Распределение
\term{spike-and-slab} задает для каждого коэффициента смесь точечной массы в
нуле и непрерывного распределения ненулевого значения; поэтому оно одновременно
описывает вероятность включения терма и неопределенность его величины
~\cite{hirsh2022}. В реализованной модели свободный член присутствует всегда,
а остальные три члена включаются независимо. Поскольку возможны только
$2^3=8$ структур,
апостериорная вероятность каждой модели вычисляется точно. Внутри модели
использовано сопряженное нормально-обратное гамма-распределение; итоговые
оценки, апостериорные вероятности включения (PIP) и интервалы усредняются по
структурам. Интервал $\kappa$ строится после нахождения устойчивого равновесия
для каждой реализации из апостериорного распределения и потому учитывает
неопределенность коэффициентов, положения равновесия и выбора членов
библиотеки.

\subsection{Neural SDE с наблюдаемым и скрытым состоянием}

В \term{Neural SDE} снос и диффузия параметризуются нейронными сетями, а
параметры оцениваются по дискретным наблюдениям через численное решение или
аппроксимацию переходов СДУ~\cite{li2020}. Базовая Neural SDE с наблюдаемым
состоянием считает каждое наблюдение истинным
состоянием и оптимизирует гауссовское правдоподобие эйлеровского перехода с
нейросетевыми $f_\theta$ и $g_\theta$. Это простой способ оценить гибкую СДУ,
но разность двух соседних зашумленных наблюдений смешивает процессный шум и
шум измерения.

\term{Neural SDE со скрытым состоянием} вводит ненаблюдаемую траекторию
$z_{0:T}$ и отдельное условное распределение измерений. Такая постановка
относится к латентным стохастическим дифференциальным моделям, в которых
динамика задается в скрытом непрерывном состоянии, а наблюдения генерируются
условно на этом состоянии~\cite{li2020}:
\begin{equation}
\label{eq:latent-sde}
 z_{t+1}\mid z_t\sim N\!\left(z_t+f_\theta(z_t)\Delta t,
 g_\theta^2(z_t)\Delta t\right),\qquad
 y_t\mid z_t\sim N(z_t,\tau^2).
\end{equation}
Параметры и факторизованное гауссовское апостериорное распределение
$q(z_{0:T})$ оптимизируются по вариационной нижней границе (ELBO). Пропуск
остается скрытым состоянием и не разрывает последовательность. Это локальный
вариационный базовый метод, а не полноценная глобальная непрерывновременная
архитектура с амортизированным энкодером; это ограничение существенно для
интерпретации.

\subsection{Модели сравнения для эпидемиологических данных}

Для классификации начала волны сравниваются Decision Tree, Random Forest, Easy
Ensemble, RUSBoost, Balanced Bagging и 25 официальных наборов предварительно
обученных весов EWSNet. EWSNet --- нейросетевой классификатор, обученный на
синтетических нормальных формах распознавать приближение различных типов
критических переходов~\cite{bury2021}. Для прогнозирования COVID-индикаторов
используются Persistence, Local Kalman-OU, сеть долгой краткосрочной памяти
(long short-term memory, LSTM)~\cite{hochreiter1997} и управляемый рекуррентный
блок (gated recurrent unit, GRU)~\cite{cho2014}. Все модели оцениваются на
одинаковых датах прогнозного отсечения. Масштабирование, балансировка и выбор
параметров выполняются без доступа к тестовому интервалу.

\section{Результаты синтетических экспериментов}

\subsection{Линейный контроль: что дает разделение механизмов}

Классическая тревога по дисперсии сработала при росте шума в 100\% траекторий
базового режима. Иными словами, она корректно заметила рост вариабельности, но
не смогла определить его причину. Простой последовательный порог по
$\widehat\kappa$ обнаружил 22.5\% траекторий потери устойчивости и дал 8.8\%
ложных тревог при росте шума; медианное упреждение среди обнаруженных составило
30 единиц времени. Средний AUC равен 0.934 для ансамбля классических EWS и
0.950 для Local SDE. H1
поддерживается в смысле специфичности, но низкая доля обнаружений показывает,
что строгая причинная тревога заметно сложнее ретроспективного различения.

Сводное сравнение точности параметров, доли успешных оценок и AUC показано на
рисунке~\ref{fig:ou-benchmark}.

\begin{figure}[htbp]
\centering
\includegraphics[width=.94\textwidth]{ou_benchmark.png}
\caption{Результаты OU-эксперимента при одинаковом горизонте}
\label{fig:ou-benchmark}
\end{figure}

На чистых данных все модели достигают AUC 1.0. Это не означает их
эквивалентности: различаются ошибки параметров и калибровка. В стрессовом
режиме Moment OU формально сохраняет AUC 0.920, но доля успешных оценок
составляет 31\% по всем траекториям и 41.7\% на тестовой выборке. Среди методов со
100\%-й долей успешных оценок Kalman OU дает AUC 0.826, а Neural SDE со скрытым
состоянием --- 0.653. Neural SDE с наблюдаемым состоянием имеет 50\%-ю долю
успешных оценок на тестовой выборке и AUC 0.400. Следовательно, H3
поддерживается прежде всего по технической устойчивости, а не по безусловному
качеству ранжирования.

Для Bayesian SINDy 95\%-й интервал $\kappa$ покрывает истинное значение в
\BayesOUBaseCoverage{} чистых окон и в \BayesOUCombinedCoverage{} окон
стрессового режима. Точное усреднение по моделям делает неопределенность
прозрачнее ARD-приближения, но не исправляет неверно заданную модель
наблюдений.

\subsection{Нелинейные переходы}

Различия между методами и нормальными формами показаны на
рисунке~\ref{fig:nonlinear-auc}.

\begin{figure}[htbp]
\centering
\includegraphics[width=.95\textwidth]{nonlinear_auc.png}
\caption{AUC по трем нормальным формам при одинаковой траектории локальной
устойчивости}
\label{fig:nonlinear-auc}
\end{figure}

В чистом нелинейном режиме лучший средний AUC дает \BestNLBaseline. При этом
ошибка физической $\kappa$ остается заметной: правильное ранжирование классов
не эквивалентно правильному восстановлению функции сноса. В стрессовом режиме
лучший формальный результат --- \BestNLCombined. Значение 0.550 близко к случайному
ранжированию и не поддерживает тезис об универсальном преимуществе гибкой
идентификации. H2 подтверждается как ограничение: дополнительная сложность
полезна лишь при достаточной локальной информации.

Эмпирическое покрытие интервалов Bayesian SINDy сопоставлено с номинальным
уровнем на рисунке~\ref{fig:bayes-calibration}.

\begin{figure}[htbp]
\centering
\includegraphics[width=.78\textwidth]{bayes_calibration.png}
\caption{Эмпирическое покрытие 95\%-го интервала Bayesian SINDy для скорости
восстановления}
\label{fig:bayes-calibration}
\end{figure}

Среднее покрытие интервалов $\kappa$ для Bayesian SINDy увеличивается с 43.3\%
в базовом режиме до 89.9\% в стрессовом. Рост покрытия возникает вместе с
расширением интервалов и не означает роста точности. Интервалы $\sigma$ в
базовом режиме имеют нулевое покрытие: модель без отдельного слоя наблюдений
приписывает шум измерения процессной диффузии. Апостериорные вероятности
включения истинных квадратичного и кубического членов также низки. На коротком
окне нелинейная нормальная форма выглядит почти линейной, а влияние шума
сильнее свидетельства в пользу глобального члена библиотеки.

\subsection{Что добавило скрытое состояние}

В нелинейном стрессовом режиме Neural SDE со скрытым состоянием имеет 100\%-ю
долю успешных оценок, тогда как эйлеровская базовая модель успешно обрабатывает
примерно половину окон. Средний AUC равен
\LatentCombinedAUC{} против \EulerCombinedAUC{}. Одновременно средний RMSE
оценки $\tau$ равен
\LatentCombinedTau{}, а AUC модели со скрытым состоянием остается около 0.5.
Значит, отдельная модель наблюдений решает часть численной проблемы, но не создает информацию,
которой нет в коротком редком ряду.

\section{Эпидемиологическая проверка}

\subsection{Критерий начала эпидемической волны}

Для каждого неэпидемического 30-недельного окна формируется положительная
метка, если ближайший официальный эпидемический период начинается не позже,
чем через $h$ недель. Окна, уже пересекающие эпидемический период, исключаются.
Модель отвечает на вопрос «начнется ли волна в ближайшие 4--5 недель?», а не
«идет ли эпидемия сейчас?». Все скользящие статистики используют
только прошлое относительно даты решения.

Лучший PR-AUC на тестовой выборке показывает Easy Ensemble: 0.713 для горизонта
4 недели и 0.747 для 5 недель. Для $h=4$ полнота равна 0.917 при точности
положительного прогноза 0.423; для $h=5$ полнота равна 1.000 при точности
0.492. Это модель раннего скрининга: она редко пропускает начало, но создает
заметное число ложных тревог.

Предварительно обученная EWSNet без дообучения присваивает почти всем окнам
переход: полнота 1.0 сочетается с PR-AUC 0.178 и 0.222. Результат поддерживает H4: перенос
синтетически обученного детектора на конкретный эпидемиологический ряд нельзя
считать успешным только из-за высокой чувствительности.

\subsection{Прогноз COVID-индикаторов со скользящим началом}

Оценивание со скользящим началом прогноза (rolling-origin evaluation)
последовательно сдвигает дату отсечения, на каждом шаге обучает или обновляет
модель только по доступному прошлому и проверяет прогноз на следующем
фиксированном горизонте~\cite{tashman2000}. Такой способ оценки создает несколько
хронологических вневыборочных проверок вместо одной произвольной границы между
обучением и тестом.

Разметка ряда ОРВИ и ошибки прогнозирования COVID-индикаторов приведены на
рисунке~\ref{fig:epidemiology}.

\begin{figure}[htbp]
\centering
\includegraphics[width=.95\textwidth]{epidemiology.png}
\caption{Слева --- размеченный недельный ряд ОРВИ; справа --- прогноз
COVID-индикаторов со скользящим началом на одинаковых датах отсечения}
\label{fig:epidemiology}
\end{figure}

На трех разбиениях с расширяющимся обучающим окном Persistence выигрывает
прогноз среднего на 14 дней, GRU --- среднего на 28 дней, а Local Kalman-OU ---
прогноз дисперсии на обоих
горизонтах. Победитель зависит от цели и горизонта; единой доминирующей
рекуррентной архитектуры нет.

\section{Обсуждение}

Рост дисперсии означает только «ряд стал более вариабельным». Падение
$\widehat\kappa$ означает «локальные отклонения затухают медленнее», но вывод
содержателен лишь при приемлемой доле успешных оценок и интервале неопределенности.
Рост $\widehat\sigma$ без падения $\widehat\kappa$ поддерживает объяснение
«усилились случайные возмущения», а не «приближается бифуркация».

Результаты различаются прежде всего по способу учета процесса наблюдения.
Методы с явной моделью наблюдений устойчивее методов, основанных на соседних
наблюдаемых приращениях. Kalman OU остается
сильной базовой моделью именно потому, что использует структурное ограничение
и отделяет $\tau$. Neural SDE со скрытым состоянием движется в том же
направлении, однако ее гибкость требует больше данных и более сильной
регуляризации.

Bayesian SINDy полезен как инструмент структурной идентификации: PIP и
интервалы, усредненные по моделям, показывают, какие члены библиотеки
поддерживаются данными и насколько неопределен выбор структуры. Однако
номинальный 95\%-й интервал нельзя автоматически считать откалиброванным.
Неверная модель наблюдений создает систематическое смещение, которое
апостериорная неопределенность внутри условной модели не покрывает.

Локальную устойчивость можно надежно восстанавливать и использовать для
различения механизмов в достаточно частых и чистых наблюдениях, особенно при
правдоподобной локальной параметризации. На коротких, редких, зашумленных и
неполных рядах задача не имеет универсального решения: доля успешных оценок падает,
нелинейные термы становятся слабо идентифицируемыми, а интервалы зависят от
описания шума измерения. Идентификация СДУ повышает интерпретируемость, но
ранняя тревога надежна лишь тогда, когда
качество данных и модель наблюдений позволяют отдельно оценить снос и
диффузию.

\textbf{Проверка гипотез.} \textbf{H1} поддержана частично: совместная оценка
$\kappa$ и $\sigma$ повышает специфичность относительно одной дисперсии, но
строгая последовательная тревога сохраняет низкую чувствительность.
\textbf{H2} поддержана: гибкие модели описывают нелинейный снос, однако на
коротких и зашумленных окнах не получают устойчивого преимущества и чаще дают
неуспешные оценки. \textbf{H3} поддержана частично: явная модель наблюдений
повышает долю успешных оценок, но не гарантирует лучшего ROC-AUC.
\textbf{H4} поддержана: EWSNet без адаптации к эпидемиологическому ряду не
сохраняет качество, полученное на синтетических данных.

Анализ ограничен одномерными вычислительными экспериментами, белым броуновским шумом,
полностью случайными пропусками (MCAR) и локальной стационарностью внутри окна.
SINDy ограничен кубической библиотекой, выбранной регуляризацией и априорным
распределением; апостериорное распределение скрытой траектории факторизовано;
переходная плотность аппроксимирована схемой Эйлера. Число нелинейных реплик
ограничено тридцатью на ячейку. Эпидемиологическая проверка носит
ретроспективный характер, поэтому выводы нельзя переносить на клинические
решения без отдельной проспективной оценки.

\section*{Заключение}
\addcontentsline{toc}{section}{Заключение}

Основные выводы по синтетическим и эпидемиологическим данным:

\begin{enumerate}
\item Рост дисперсии не идентифицирует механизм изменения. В линейном
эксперименте тревога по дисперсии сработала при одном лишь росте процессного
шума в 100\% траекторий базового режима.
\item Совместная оценка скорости восстановления $\kappa$ и диффузии $\sigma$
повышает специфичность вывода. Local SDE достиг AUC 0.950, но строгая
последовательная тревога обнаружила только 22.5\% переходов при 8.8\% ложных
тревог в сценарии роста шума.
\item Явная модель наблюдений повышает вычислительную устойчивость при редких
и неполных данных. В стрессовом OU-режиме Kalman OU вернул допустимую оценку во
всех тестовых окнах; модель Neural SDE со скрытым состоянием также была
устойчивее варианта, отождествляющего наблюдение с истинным состоянием.
\item Гибкая нелинейная идентификация не получила универсального преимущества.
Лучший средний AUC в нелинейном стрессовом режиме равен 0.550, поэтому данные
не подтверждают надежное различение механизмов на текущей длине окна.
\item Эпидемиологическая проверка воспроизводит критерий начала волны. Easy
Ensemble достиг PR-AUC 0.713 и 0.747 на горизонтах 4 и 5 недель; для
COVID-индикаторов победитель зависел от цели и горизонта прогноза.
\end{enumerate}

Итог проверки гипотез следующий: \textbf{H1} и \textbf{H3} получили частичную
поддержку, поскольку специфичность и вычислительная устойчивость улучшились без
универсального роста чувствительности или ROC-AUC; \textbf{H2} и \textbf{H4}
поддержаны результатами нелинейных и эпидемиологических экспериментов.

Дальнейший анализ включает следующие этапы:

\begin{enumerate}
\item до дальнейшего расширения моделей зафиксировать подтверждающее разбиение,
горизонты, основную метрику и порог тревоги;
\item добавить анализ чувствительности к априорному распределению и длине окна,
бутстреп-интервалы для разностей метрик и апостериорные предиктивные проверки;
\item дополнить модель со скрытым состоянием амортизированным энкодером и
проверить частичную наблюдаемость многомерных систем;
\item протестировать цветной и скачкообразный шум, неслучайные пропуски и изменение
частоты наблюдений;
\item после фиксации правил сравнения повторить анализ на независимом
эпидемиологическом наборе и отдельно оценить калибровку вероятности тревоги.
\end{enumerate}

Предварительная фиксация правил анализа снижает риск подстройки к данным.

\clearpage
\section*{Использование искусственного интеллекта}
\addcontentsline{toc}{section}{Использование искусственного интеллекта}

При подготовке работы генеративные инструменты использовались для разработки и
рефакторинга программного кода, подготовки репозитория, технической редактуры,
унификации терминологии и набора отдельных фрагментов в LaTeX. Численные
результаты воспроизводятся кодом из репозитория на зафиксированных данных и
параметрах. Ответственность
за окончательную проверку расчетов, интерпретацию результатов и представленный
текст несет исключительно автор работы.

\clearpage

\renewcommand{\refname}{\centering СПИСОК ИСПОЛЬЗОВАННЫХ ИСТОЧНИКОВ}
\begin{thebibliography}{99}
\addcontentsline{toc}{section}{Список использованных источников}
\sloppy

\bibitem{scheffer2009}
Scheffer M. Early-warning signals for critical transitions / M. Scheffer,
J. Bascompte, W. A. Brock [et al.] // Nature. --- 2009. --- Vol. 461. ---
P. 53--59. --- DOI 10.1038/nature08227.

\bibitem{bonalli2026}
Bonalli R. Non-parametric learning of stochastic differential equations with
non-asymptotic fast rates of convergence / R. Bonalli, A. Rudi // Foundations
of Computational Mathematics. --- 2026. --- Vol. 26. --- P. 1497--1552. ---
DOI 10.1007/s10208-025-09705-x.

\bibitem{strogatz2015}
Strogatz S. H. Nonlinear dynamics and chaos: with applications to physics,
biology, chemistry, and engineering. --- 2nd ed. --- Boulder : Westview Press,
2015. --- 513 p.

\bibitem{kalman1960}
Kalman R. E. A new approach to linear filtering and prediction problems //
Journal of Basic Engineering. --- 1960. --- Vol. 82, no. 1. --- P. 35--45. ---
DOI 10.1115/1.3662552.

\bibitem{carpenter2011}
Carpenter S. R. Early warnings of unknown nonlinear shifts: a nonparametric
approach / S. R. Carpenter, W. A. Brock // Ecology. --- 2011. --- Vol. 92,
no. 12. --- P. 2196--2201. --- DOI 10.1890/11-0716.1.

\bibitem{lade2012}
Lade S. J. Early warning signals for critical transitions: a generalized
modeling approach / S. J. Lade, T. Gross // PLoS Computational Biology. ---
2012. --- Vol. 8, no. 2. --- Art. e1002360. --- DOI
10.1371/journal.pcbi.1002360.

\bibitem{bury2021}
Bury T. M. Deep learning for early warning signals of tipping points /
T. M. Bury, R. I. Sujith, I. Pavithran [et al.] // Proceedings of the National
Academy of Sciences. --- 2021. --- Vol. 118, no. 39. --- Art. e2106140118. ---
DOI 10.1073/pnas.2106140118.

\bibitem{rubin1976}
Rubin D. B. Inference and missing data // Biometrika. --- 1976. --- Vol. 63,
no. 3. --- P. 581--592. --- DOI 10.1093/biomet/63.3.581.

\bibitem{lazutov2026}
Лазутов М. Ю. Прогнозирование эпидемий острых респираторных инфекций с
использованием сигналов раннего предупреждения и моделей машинного обучения /
М. Ю. Лазутов, В. Н. Леоненко // Автоматика и телемеханика. --- 2026. ---
№ 6. --- DOI 10.7868/S2413977726060056.

\bibitem{saito2015}
Saito T. The precision-recall plot is more informative than the ROC plot when
evaluating binary classifiers on imbalanced datasets / T. Saito,
M. Rehmsmeier // PLoS ONE. --- 2015. --- Vol. 10, no. 3. --- Art. e0118432. ---
DOI 10.1371/journal.pone.0118432.

\bibitem{brunton2016}
Brunton S. L. Discovering governing equations from data by sparse
identification of nonlinear dynamical systems / S. L. Brunton, J. L. Proctor,
J. N. Kutz // Proceedings of the National Academy of Sciences. --- 2016. ---
Vol. 113, no. 15. --- P. 3932--3937. --- DOI 10.1073/pnas.1517384113.

\bibitem{hirsh2022}
Hirsh S. M. Sparsifying priors for Bayesian uncertainty quantification in model
discovery / S. M. Hirsh, D. A. Barajas-Solano, J. N. Kutz // Royal Society Open
Science. --- 2022. --- Vol. 9, no. 2. --- Art. 211823. --- DOI
10.1098/rsos.211823.

\bibitem{li2020}
Li X. Scalable gradients for stochastic differential equations / X. Li,
T.-K. L. Wong, R. T. Q. Chen, D. Duvenaud // Proceedings of Machine Learning
Research. --- 2020. --- Vol. 108. --- P. 3870--3882.

\bibitem{hochreiter1997}
Hochreiter S. Long short-term memory / S. Hochreiter, J. Schmidhuber // Neural
Computation. --- 1997. --- Vol. 9, no. 8. --- P. 1735--1780. --- DOI
10.1162/neco.1997.9.8.1735.

\bibitem{cho2014}
Cho K. Learning phrase representations using RNN encoder--decoder for
statistical machine translation / K. Cho, B. van Merriënboer, C. Gulcehre
[et al.] // Proceedings of the 2014 Conference on Empirical Methods in Natural
Language Processing. --- Doha : Association for Computational Linguistics,
2014. --- P. 1724--1734. --- DOI 10.3115/v1/D14-1179.

\bibitem{tashman2000}
Tashman L. J. Out-of-sample tests of forecasting accuracy: an analysis and
review // International Journal of Forecasting. --- 2000. --- Vol. 16, no. 4.
--- P. 437--450. --- DOI 10.1016/S0169-2070(00)00065-0.

\end{thebibliography}

\clearpage
\appendix
\makeatletter
\@addtoreset{table}{section}
\makeatother
\renewcommand{\thetable}{\Alph{section}.\arabic{table}}
\section{Полные таблицы результатов}

\subsection{Сравнительный эксперимент с OU-процессом}

Полные результаты линейного эксперимента приведены в
таблицах~\ref{tab:ou-summary} и~\ref{tab:ou-alarms}.

\begin{table}[H]
\centering\small
\caption{Восстановление параметров на отложенной выборке и различение потери
устойчивости и роста шума}
\label{tab:ou-summary}
\resizebox{\textwidth}{!}{\input{report/tables/ou_summary.tex}}
\end{table}

\begin{table}[H]
\centering\small
\caption{Последовательные тревоги OU-процесса после калибровки по траекториям}
\label{tab:ou-alarms}
\resizebox{\textwidth}{!}{\input{report/tables/ou_alarms.tex}}
\end{table}

\subsection{Нелинейные нормальные формы}

Результаты по нелинейным системам детализированы в
таблицах~\ref{tab:nl-summary}--\ref{tab:nl-alarms}.

\begin{table}[H]
\centering\small
\caption{Средние по трем нормальным формам метрики на отложенной выборке}
\label{tab:nl-summary}
\resizebox{\textwidth}{!}{\input{report/tables/nonlinear_summary.tex}}
\end{table}

\begin{table}[H]
\centering\small
\caption{Эмпирическая калибровка spike-and-slab на нелинейных системах}
\label{tab:bayes-calibration}
\resizebox{\textwidth}{!}{\input{report/tables/bayes_calibration.tex}}
\end{table}

\begin{table}[H]
\centering\small
\caption{Средние апостериорные вероятности включения в поздних окнах сценария
потери устойчивости}
\label{tab:bayes-pip}
\input{report/tables/bayes_pip.tex}
\end{table}

\begin{table}[H]
\centering\small
\caption{Последовательные тревоги в нелинейных системах с усреднением по
нормальным формам}
\label{tab:nl-alarms}
\resizebox{\textwidth}{!}{\input{report/tables/nonlinear_alarms.tex}}
\end{table}

\subsection{Эпидемиологические результаты}

Полные метрики эпидемиологических экспериментов приведены в
таблицах~\ref{tab:orvi-results} и~\ref{tab:covid-winners}.

\begin{table}[H]
\centering\small
\caption{Тестовые метрики классификаторов начала волны ОРВИ}
\label{tab:orvi-results}
\resizebox{\textwidth}{!}{\input{report/tables/orvi_results.tex}}
\end{table}

\begin{table}[H]
\centering\small
\caption{Победители по среднему NRMSE в схеме со скользящим началом прогноза}
\label{tab:covid-winners}
\input{report/tables/covid_rolling_winners.tex}
\end{table}

\end{document}
